% !TeX root = ../../book.tex
\section{逻辑等价}
\secbegin{secLogicalEquivalence}

我们用一个例子来引出本节的内容。

\begin{example}
\label{exNoEvenPrimeGreaterThanTwo}
考虑以下两个逻辑公式，其中 $P$ 表示所有质数的集合。
\begin{enumerate}[(1)]
\item $\forall n \in P,\, (n > 2 \Rightarrow [\exists k \in \mathbb{Z},\, n=2k+1])$;
\item $\neg \exists n \in P,\, (n > 2 \wedge [\exists k \in \mathbb{Z},\, n=2k])$.
\end{enumerate}
逻辑公式（1）可以翻译为`每个大于二的质数都是奇数'，逻辑公式（2）翻译为`不存在大于二的偶质数'。这两个陈述显然是 \textit{等价的} --- 它们的意思相同 --- 但它们提出了不同的证明策略：

\begin{enumerate}[(1)]
\item 固定一个质数 $n$，假设 $n>2$，然后证明对于某些 $k \in \mathbb{Z}$, $n = 2k+1$ 。
\item 假设存在某个质数 $n$ 使得 $n>2$ 且对于某些 $k \in \mathbb{Z}$, $n=2k$，并推导出矛盾。
\end{enumerate}

虽然陈述 (1) 更直接地翻译成中文陈述`每个大于二的素数都是奇数'，但陈述 (2) 所建议的证明策略更易于使用。实际上，如果 $n$ 是质数，使得 $n>2$ 且对于某些 $k \in \mathbb{Z}$, $n=2k$，则 $2$ 是 $n$ 的约数，不只有 $1$ 和 $n$（因为 $1<2<n$），与 $n$ 是质数的假设相矛盾。
\end{example}

\textit{逻辑等价}的概念准确地抓住了 \Cref{exNoEvenPrimeGreaterThanTwo} 中 (1) 和 (2) 逻辑公式`意思相同'的意义。能够将逻辑公式转换为不同（但等价）的形式使我们能够确定更广泛的可行证明策略。

\begin{definition}
\label{defLogicalEquivalence}
\index{logical equivalence}
\index{equivalence!logical}
\nindex{logical equivalence}{$\equiv$}{logical equivalence}
设 $p$ 和 $q$ 为逻辑公式。如果 $q$ 可以从 $p$ 推导出，$p$ 可以从 $q$ 推导出，则我们说 $p$ 和 $q$ 是\textbf{逻辑等价}的，写做 $p \equiv q$ \inlatex{equiv}\lindexmmc{equiv}{$\equiv$}。
\end{definition}

\subsection*{命题公式的逻辑等价}

虽然 \Cref{defLogicalEquivalence} 定义了任意逻辑公式之间的逻辑等价，但我们将首先关注\textit{命题}公式之间的逻辑等价，就像我们在 \Cref{secPropositionalLogic} 中看到的那样。

首先，我们先看几个例子，看看逻辑等价的证明可能是什么样子。前方高能 --- 他们读起来不是很友好！但隧道的尽头会有光。在经历了 \Cref{exConjunctionDistributesOverDisjunction,exImplicationInTermsOfDisjunction} 和 \Cref{exPAndQImpliesRIffPImpliesRAndQImpliesR} 之后，我们将介绍一个非常快速而简单的工具来证明命题公式在逻辑上是等价的。

\begin{example}
\label{exConjunctionDistributesOverDisjunction}
证明 $p \wedge (q \vee r) \equiv (p \wedge q) \vee (p \wedge r)$，其中 $p$、$q$ 和 $r$ 是命题变量。

\begin{itemize}
\item 首先假设 $p \wedge (q \vee r)$ 为真。那么根据合取的定义，$p$ 为真且 $q \vee r$ 为真。根据析取的定义，$q$ 为真或 $r$ 为真。
\begin{itemize}
\item 如果 $q$ 为真，则 $p \wedge q$ 根据合取的定义为真。
\item 如果 $r$ 为真，则 $p \wedge r$ 根据合取的定义为真。
\end{itemize}
在这两种情况下，根据析取的定义，$(p \wedge q) \vee (p \wedge r)$ 为真。

\item 假设 $(p \wedge q) \vee (p \wedge r)$ 为真。那么根据析取的定义，要么 $p \wedge q$ 为真，要么 $p \wedge r$ 为真。
\begin{itemize}
\item 如果 $p \wedge q$ 为真，则根据合取的定义，$p$ 为真且 $q$ 为真。
\item 如果 $p \wedge r$ 为真，则根据合取的定义，$p$ 为真且 $r$ 为真。
\end{itemize}
在这两种情况下，根据析取的定义，$p$ 为真，$q \vee r$ 为真。因此，根据合取的定义，$p \wedge (q \vee r)$ 为真。
\end{itemize}

因为我们可以从 $p \wedge (q \vee r)$ 推导出 $(p \wedge q) \vee (p \wedge r)$，反之亦然，它满足
\[ p \wedge (q \vee r) \equiv (p \wedge q) \vee (p \wedge r)\]
的要求。
\end{example}

\begin{example}
\label{exImplicationInTermsOfDisjunction}
我们证明 $p \Rightarrow q \equiv (\neg p) \vee q$，其中 $p$、$q$ 和 $r$ 是命题变量。

\begin{itemize}
\item 首先假设 $p \Rightarrow q$ 为真。根据排中律 (\Cref{axLE})，要么 $p$ 为真，要么 $\neg p$ 为真 --- 在每种情况下我们都导出 $(\neg p) \vee q$。
\begin{itemize}
\item 如果 $p$ 为真，则因为 $p \Rightarrow q$ 为真, 由 \elimrule{\Rightarrow} 可得 $q$ 为真，因此由 \introrulesub{\vee}{2} 可得 $(\neg p) \vee q$ 为真；
\item 如果 $\neg p$ 为真，则根据 \introrulesub{\vee}{1} 可得 $(\neg p) \vee q$ 为真。
\end{itemize}
两种情况下，我们得到 $(\neg p) \vee q$ 为真。

\item 假设 $(\neg p) \vee q$ 为真。证明 $p \Rightarrow q$ 为真，根据 \introrule{\Rightarrow} ，假设 $p$ 为真并推导出 $q$ 就可以了。所以假设 $p$ 为真。由于 $(\neg p) \vee q$ 为真，所以我们有 $\neg p$ 为真或 $q$ 为真。
\begin{itemize}
\item 如果 $\neg p$ 为真，则与假设 $p$ 为真矛盾，因此根据爆炸原理 (\Cref{axPrincipleOfExplosion}) $q$ 为真。
\item 如果 $q$ 为真\dots{} 那么 $q$ 就是真的 --- 没有什么要证明的了！
\end{itemize}
两种情况下，我们都得到 $q$ 为真。因此 $p \Rightarrow q$ 为真。
\end{itemize}

我们从 $p \Rightarrow q$ 推导出 $(\neg p) \vee q$ ，反之亦然，因此这两个公式逻辑等价。
\end{example}

\begin{exercise}
\label{exPAndQImpliesRIffPImpliesRAndQImpliesR}
设 $p$、$q$ 和 $r$ 为命题变量。证明命题公式$(p \vee q) \Rightarrow r$ 在逻辑上等价于 $(p \Rightarrow r) \wedge (q \Rightarrow r)$。
\end{exercise}

证明逻辑等价的推导可能会很麻烦，即使是对于像 \Cref{exConjunctionDistributesOverDisjunction,exImplicationInTermsOfDisjunction} 和 \Cref{exPAndQImpliesRIffPImpliesRAndQImpliesR} 这样的小例子也是如此。

下面的定理将证明\textit{命题}公式之间的逻辑等价性问题简化为纯粹的算法任务，即在（相对）小的情况列表中检查公式何时为真，何时为假。我们将使用 \textit{真值表} (\Cref{defTruthTable}) 进一步简化这个过程。

\begin{theorem}
\label{thmLogicalEquivalentIffSameTruthValues}
两个命题公式逻辑等价，当且仅当构成命题的变量取任何真值下，两个命题的真值都是相同的。
\end{theorem}

\begin{cidea}
尽管我们可以通过 \Cref{secStructuralInduction} 中介绍的\textit{结构归纳}严格证明上面的定理，但这在当下有点超纲。

证明思路是，由于命题公式是使用逻辑运算符从更简单的命题公式构建的，因此更复杂的命题公式的真值由其更简单的子公式的真值决定。如果我们继续`追踪'这些子公式，我们最终只会得到命题变量。

例如 $(p \Rightarrow r) \wedge (q \Rightarrow r)$ 的真值由 $p \Rightarrow r$ 和 $q \Rightarrow r$ 的真值根据合取运算符 $\wedge$ 的规则确定 。接着，$p \Rightarrow r$ 的真值由 $p$ 和 $r$ 的真值根据蕴涵运算符 $\Rightarrow$ 决定，$q \Rightarrow r$ 的真值由 $q$ 和 $r$ 的真值再次根据蕴涵运算符决定。由此可见，整个命题公式 $(p \Rightarrow r) \wedge (q \Rightarrow r)$ 的真值根据 $\wedge$ 和 $\Rightarrow$ 的规则由 $p,q,r$ 的真值决定。

如果将真值赋给命题变量使得一个命题公式为真而另一个为假，那么肯定不可能从一个推导出另一个 --- 否则我们会得到一个矛盾。因此，无论对命题变量的真值赋什么值，两个命题公式必须具有相同的真值。
\end{cidea}

现在，我们开发一套系统的方法来检查命题变量赋值后命题公式的真值。

\begin{definition}
\label{defTruthTable}
\index{truth table}
命题公式的\textbf{真值表}是一张表，其中每行表示命题变量可能的赋值，每列表示一个子公式（包括命题变量和命题公式本身）。真值表的条目是子公式的真值。
\end{definition}

\begin{example}
\label{exNegationConjunctionDisjunctionImplicationTruthTable}
一下是 $\neg p$, $p \wedge q$, $p \vee q$ 和 $p \Rightarrow q$ 的真值表。

\begin{center}
\begin{tabular}{c|c}
$p$ & $\neg p$ \\ \hline
\TT & \FF \\
\FF & \TT \\
\multicolumn{2}{c}{\phantom{\TT}\phantom{\FF}} \\
\multicolumn{2}{c}{\phantom{\TT}\phantom{\FF}}
\end{tabular}
%
\hspace{15pt}
%
\begin{tabular}{cc|c}
$p$ & $q$ & $p \wedge q$ \\ \hline
\TT & \TT & \TT \\
\TT & \FF & \FF \\
\FF & \TT & \FF \\
\FF & \FF & \FF
\end{tabular}
%
\hspace{15pt}
%
\begin{tabular}{cc|c}
$p$ & $q$ & $p \vee q$ \\ \hline
\TT & \TT & \TT \\
\TT & \FF & \TT \\
\FF & \TT & \TT \\
\FF & \FF & \FF
\end{tabular}
%
\hspace{15pt}
%
\begin{tabular}{cc|c}
$p$ & $q$ & $p \Rightarrow q$ \\ \hline
\TT & \TT & \TT \\
\TT & \FF & \FF \\
\FF & \TT & \TT \\
\FF & \FF & \TT
\end{tabular}
\end{center}
\end{example}

在 \Cref{exNegationConjunctionDisjunctionImplicationTruthTable} 中，我们使用符号 \TT{} \inlatex{checkmark}\lindexmmc{checkmark}{$\checkmark$} 表示`真'，用 \FF{} \inlatex{times}\lindexmmc{times }{$\times$} 表示`假'。一些作者采用其他约定，例如 $T,F$ 或 $\top,\bot$ \inlatex{top,\textbackslash{}bot}\lindexmmc{top}{$\top$}\lindexmmc{bot}{$ \bot$} 或 $1,0$ 或 $0,1$--- 无限可能！

\begin{exercise}
使用 $\wedge$、$\vee$ 和 $\Rightarrow$ 的定义来证明 \Cref{exNegationConjunctionDisjunctionImplicationTruthTable} 中的真值表。
\hintlabel{exJustifyBasicTruthTables}{请注意，您可能需要使用排中律 (\Cref{axLEM}) 和爆炸原理 (\Cref{axPrincipleOfExplosion})。}
\end{exercise}

下一个例子展示了如何将各个逻辑运算符的真值表（如 \Cref{exNegationConjunctionDisjunctionImplicationTruthTable}）组合起来，形成一个涉及三个命题变量的复杂命题公式的真值表。

\begin{example}
\label{exFirstExampleOfTruthTable}
以下是 $(p \wedge q) \vee (p \wedge r)$ 的真值表。
\begin{center}
\begin{tabular}{ccc|cc|c}
$p$ & $q$ & $r$ & $p \wedge q$  & $p \wedge r$ & $(p \wedge q) \vee (p \wedge r)$\\ \hline
\TT & \TT & \TT & \TT           & \TT          & \TT \\
\TT & \TT & \FF & \TT           & \FF          & \TT \\
\TT & \FF & \TT & \FF           & \TT          & \TT \\
\TT & \FF & \FF & \FF           & \FF          & \FF \\
\FF & \TT & \TT & \FF           & \FF          & \FF \\
\FF & \TT & \FF & \FF           & \FF          & \FF \\
\FF & \FF & \TT & \FF           & \FF          & \FF \\
\FF & \FF & \FF & \FF           & \FF          & \FF \\
\multicolumn{3}{c}{\upbracefill} & \multicolumn{2}{c}{\upbracefill} & \multicolumn{1}{c}{\upbracefill} \\
\multicolumn{3}{c}{\begin{minipage}{40pt}\centering\scriptsize 命题变量\end{minipage}} & \multicolumn{2}{c}{\begin{minipage}{40pt}\centering\scriptsize 中间子公式\end{minipage}} & \multicolumn{1}{c}{\scriptsize 主公式}
\end{tabular}
\end{center}
下面是关于这个真值表构造的一些详细解释：
\begin{itemize}
\item 首先列出命题变量。因为有 3 个，所以有 $2^3=8$ 行。$p$ 列包含四个 \TT{} 后跟四个 \FF{}; $q$ 列包含两个 \TT{}，两个 \FF{}，然后重复；$r$ 列包含一个 \TT{}，一个 \FF{}，然后重复。
\item 下一组列是较复杂的子公式。每个子公式都是通过查看左侧相关列并与合取真值表进行比较来构建的。
\item 最后一列是主公式本身，它也是通过查看左侧相关列并与析取真值表进行比较来构造的。
\end{itemize}
竖线的放置位置以及行的放置顺序并不是唯一的，但是在为更复杂的逻辑公式构建真值表时，开发一个系统并坚持下去很有必要。
\end{example}

回到\Cref{thmLogicalEquivalentIffSameTruthValues}，我们得到以下证明两个命题公式逻辑等价的策略。

\begin{strategy}[用真值表证明逻辑等价]
\label{strTruthTable}
要想证明命题公式逻辑等价，只要证明它们有相同的真值表即可。
\end{strategy}

\begin{example}
在\Cref{exConjunctionDistributesOverDisjunction}中我们证明了 $p \wedge (q \vee r) \equiv (p \wedge q) \vee (p \wedge r)$。我们使用真值表再证明一次。首先构造$p \wedge (q \vee r)$的真值表：
\begin{center}
\begin{tabular}{ccc|c|c}
$p$ & $q$ & $r$ & $q \vee r$ & $p \wedge (q \vee r)$ \\ \hline
\TT & \TT & \TT & \TT & \TT \\
\TT & \TT & \FF & \TT & \TT \\
\TT & \FF & \TT & \TT & \TT \\
\TT & \FF & \FF & \FF & \FF \\
\FF & \TT & \TT & \TT & \FF \\
\FF & \TT & \FF & \TT & \FF \\
\FF & \FF & \TT & \TT & \FF \\
\FF & \FF & \FF & \FF & \FF
\end{tabular}
\end{center}
请注意，$p \wedge (q \vee r)$ 列与 \Cref{exFirstExampleOfTruthTable} 中 $(p \wedge q) \vee (p \wedge r)$ 列的真值相同。 因此这两个公式逻辑等价。
\end{example}

为了避免写两个真值表，将它们合并成一个真值表。例如，下面的真值表表明 $p \wedge (q \vee r)$ 在逻辑上等价于 $(p \wedge q) \vee (p \wedge r)$：

\begin{center}
\begin{tabular}{ccc||c|c||cc|c}
$p$ & $q$ & $r$ & $q \vee r$ & $p \wedge (q \vee r)$ & $p \wedge q$  & $p \wedge r$ & $(p \wedge q) \vee (p \wedge r)$\\ \hline
\TT & \TT & \TT & \TT & \TT & \TT           & \TT          & \TT \\
\TT & \TT & \FF & \TT & \TT & \TT           & \FF          & \TT \\
\TT & \FF & \TT & \TT & \TT & \FF           & \TT          & \TT \\
\TT & \FF & \FF & \FF & \FF & \FF           & \FF          & \FF \\
\FF & \TT & \TT & \TT & \FF & \FF           & \FF          & \FF \\
\FF & \TT & \FF & \TT & \FF & \FF           & \FF          & \FF \\
\FF & \FF & \TT & \TT & \FF & \FF           & \FF          & \FF \\
\FF & \FF & \FF & \FF & \FF & \FF           & \FF          & \FF
\end{tabular}
\end{center}

在下面的练习中，我们使用真值表来重写 \Cref{exImplicationInTermsOfDisjunction} 和 \Cref{exPAndQImpliesRIffPImpliesRAndQImpliesR} 的逻辑等价证明。

\begin{exercise}
\label{exImplicationInTermsOfDisjunctionWithTruthTables}
用真值表证明 $p \Rightarrow q \equiv (\neg p) \vee q$。
\end{exercise}

\begin{exercise}
\label{exPAndQImpliesRIffPImpliesRAndQImpliesRWithTruthTables}
设 $p$、$q$ 和 $r$ 为命题变量。用真值表证明命题公式 $(p \vee q) \Rightarrow r$ 在逻辑上等价于 $(p \Rightarrow r) \wedge(q \Rightarrow r)$。
\end{exercise}

\subsection*{一些证明策略}

现在，我们可以很好地使用逻辑等价来推导一些更复杂的证明策略。

\begin{theorem}[双重否定律]
\label{thmDoubleNegation}
设 $p$ 为命题变量，则 $p \equiv \neg \neg p$。
\end{theorem}

\begin{cproof}
使用真值表可以轻而易举地证明出来。事实上，请考虑以下真值表。
\begin{center}
\begin{tabular}{c|c|c}
$p$ & $\neg p$ & $\neg \neg p$ \\ \hline
\TT & \FF & \TT \\
\FF & \TT & \FF
\end{tabular}
\end{center}
$p$ 列和 $\neg \neg p$ 列是相同的，因此 $p \equiv \neg \neg p$。
\end{cproof}

双重否定律很重要，因为它给出了通过矛盾来证明陈述的第二种方法。事实上，它说的是证明一个命题 $p$ 相当于证明 $\neg \neg p$，这相当于假设 $\neg p$ 并导出矛盾。

\begin{strategy}[反证法（间接版）]
\label{strProofByContradictionIndirect}
\index{contradiction!(indirect) proof by}
\index{proof!by contradiction (indirect)}
为了证明命题 $p$ 为真，只需假设 $p$ 为假并导出矛盾即可。
\end{strategy}

乍一看，\Cref{strProofByContradictionIndirect} 看起来与 \Cref{strProvingNegationsDirect} 非常相似，我们也将其称为 \textit{矛盾证明}。但两者之间有一个重要的区别：
\begin{itemize}
\item \Cref{strProvingNegationsDirect} 说的是，要想证明一个命题为\textit{假}，只需假设它为\textit{真}并导出矛盾即可；
\item \Cref{strProofByContradictionIndirect} 说的是，要想证明一个命题为\textit{真}，只需假设它为\textit{假} 并导出矛盾即可。
\end{itemize}

前者是一种\textit{直接}证明技术，因为它直接源于否定运算符的定义；后者是一种\textit{间接}证明技术，因为它源于逻辑等价，即双重否定律。

\begin{example}
我们要证明，如果 $a$、$b$ 和 $c$ 是满足 $a^2+b^2=c^2$ 的非负实数，则 $a+b \ge c$。

事实上，令 $a,b,c \in \mathbb{R}$ 且 $a,b,c \ge 0$，并假设 $a^2+b^2=c^2$。为了引出矛盾，假设 $a+b \ge c$ 不成立。则必有 $a+b < c$。那么
\[ (a+b)^2 = (a+b)(a+b) < (a+b) c < c \cdot c = c^2\]
因此
\[ c^2 > (a+b)^2 = a^2+2ab+b^2 = c^2+2ab \ge c^2\]
这意味着 $c^2 > c^2$，这明显矛盾。因此按要求一定是 $a + b \ge c$。
\end{example}

我们得出的下一个证明策略涉及证明蕴涵。

\begin{definition}
\label{defContrapositive}
\index{contrapositive}
$p \Rightarrow q$ 形式的命题的\textbf{逆否}命题是 $\neg q \Rightarrow \neg p$。
\end{definition}

\begin{theorem}[逆反律]
\label{thmLawOfContraposition}
设 $p$ 和 $q$ 为命题变量。则 $p \Rightarrow q \equiv (\neg q) \Rightarrow (\neg p)$。
\end{theorem}

\begin{cproof}
我们构建 $p \Rightarrow q$ 和 $(\neg q) \Rightarrow (\neg p)$ 的真值表。

\begin{center}
\begin{tabular}{cc||c||cc|c}
$p$ & $q$ & $p \Rightarrow q$ & $\neg q$ & $\neg p$ & $(\neg q) \Rightarrow (\neg p)$ \\ \hline
\TT & \TT & \TT & \FF & \FF & \TT \\
\TT & \FF & \FF & \TT & \FF & \FF \\
\FF & \TT & \TT & \FF & \TT & \TT \\
\FF & \FF & \TT & \TT & \TT & \TT
\end{tabular}
\end{center}

$p \Rightarrow q$ 列和 $(\neg q) \Rightarrow (\neg p)$ 列相同，因此它们在逻辑上是等价的。
\end{cproof}

\Cref{thmLawOfContraposition} 给出以下证明策略。

\begin{strategy}[对位证明法]
\label{strProofByContraposition}
\index{contraposition!proof by}
\index{proof!by contraposition}
为了证明 $p \Rightarrow q$ 形式的命题，只需假设 $q$ 为假并推导出 $p$ 为假。
\end{strategy}

\begin{example}
给定两个自然数 $m$ 和 $n$。证明，如果 $mn > 64$，则 $m>8$ 或 $n>8$。

根据对位法，只需假设 $m > 8$ 或 $n > 8$ 不成立，并推导出 $mn > 64$ 不成立。

因此，假设 $m>8$ 和 $n>8$ 都不成立。则 $m \le 8$ 且 $n \le 8$，因此 $mn \le 64$。
\end{example}

\begin{exercise}
使用逆反律证明 $p \Leftrightarrow q \equiv (p \Rightarrow q) \wedge ((\neg p) \Rightarrow (\neg q))$，并使用该等价给出的证明技术来证明一个整数是偶数当且仅当它的平方是偶数。
\hintlabel{exIntegerEvenIffSquareEven}{%
将此语句表示为 $\forall n \in \mathbb{Z},\, (n \text{为偶数}) \Leftrightarrow (n^2 \text{为偶数})$，并注意 `$x$ 为偶数' 的逆命题是 `$x$ 为奇数'。
}
\end{exercise}

引用听起来令人印象深刻的结果（例如\textit{对位证明法}）感觉很好，但在实践中，\textit{任何}两个不同命题公式之间的逻辑等价都暗示着一种新的证明技术，而且并非所有这些技术都有名称。事实上，下面练习中的证明策略虽然有用，但却没有一个听起来很华而不实的名字 --- 至少没有一个能被广泛理解的名字。

\begin{exercise}
证明 $p \vee q \equiv (\neg p) \Rightarrow q$。使用该逻辑等价给出证明 $p \vee q$ 形式命题的新的证明策略，并使用该策略证明如果两个整数之和为偶数，则两个整数要么都是偶数，要么都是奇数。
\end{exercise}

\subsection*{否定}

在纯数学中，经常会问数学对象的某个属性是否成立。例如，在 \Cref{secCompletenessConvergence} 中，我们将研究实数序列的收敛性：假设实数序列 $x_0,x_1,x_2,\dots$ \textit{收敛}\label{pConvergencePreliminary} (\Cref{defConvergenceOfSequence }) 就是说
\[ \exists a \in \mathbb{R},\, \forall \varepsilon \in \mathbb{R},\, (\varepsilon > 0 \Rightarrow \exists N \in \mathbb{N},\, \forall n \in \mathbb{N},\, [n \ge N \Rightarrow |x_n-a| < \varepsilon])\]
这已经是一个比较复杂的逻辑公式了。但是如果我们想证明一个序列\textit{不}收敛怎么办？简单地假设上面的逻辑公式成立并推导出矛盾有时可能会起作用，但并不是特别有启发性。

我们的下一目标是开发一种系统方法来否定复杂的逻辑公式。完成此操作后，我们将能够否定表示`序列 $x_0, x_1, x_2, \dots$ 收敛'的逻辑公式，如下所示
\[ \forall a \in \mathbb{R},\, \exists \varepsilon \in \mathbb{R},\, (\varepsilon > 0 \wedge \forall N \in \mathbb{N},\, \exists n \in \mathbb{N},\, [n \ge N \wedge |x_n-a| \ge \varepsilon])\]

诚然，这仍然是一个复杂的表达式，但是当逐个元素分解后，它提供了有关如何证明它的有用信息。

否定合取和析取法则是\textit{德摩根定律}的实例，它揭露出 $\wedge$ 和 $\vee$ 之间的一种对偶性。

\begin{theorem}[逻辑运算符德摩根定律]
\label{thmDeMorganLogicalOperators}
\index{de Morgan's laws!for logical operators}
设 $p$ 和 $q$ 为逻辑公式。则：
\begin{enumerate}[(a)]
\item $\neg (p \wedge q) \equiv (\neg p) \vee (\neg q)$；且
\item $\neg (p \vee q) \equiv (\neg p) \wedge (\neg q)$。
\end{enumerate}
\end{theorem}

\begin{cproof}[of (a)]
考虑下面的真值表。
\begin{center}
\begin{tabular}{cc||c|c||cc|c}
$p$ & $q$ & $p \wedge q$ & $\neg (p \wedge q)$ & $\neg p$ & $\neg q$ & $(\neg p) \vee (\neg q)$ \\ \hline
\TT & \TT & \TT & \FF & \FF & \FF & \FF \\
\TT & \FF & \FF & \TT & \FF & \TT & \TT \\
\FF & \TT & \FF & \TT & \TT & \FF & \TT \\
\FF & \FF & \FF & \TT & \TT & \TT & \TT \\
\end{tabular}
\end{center}
$\neg (p \wedge q)$ 列和 $(\neg p) \vee (\neg q)$ 列相同，因此它们在逻辑上是等价的。
\end{cproof}

\begin{exercise}
用三种方法证明 \Cref{thmDeMorganLogicalOperators}(b)：一次直接使用逻辑等价的定义（就像我们在 \Cref{exConjunctionDistributesOverDisjunction,exImplicationInTermsOfDisjunction} 和 \Cref{exPAndQImpliesRIffPImpliesRAndQImpliesR} 中所做的那样），一次使用真值表，一次使用 (a) 和双重否定定律。
\end{exercise}

\begin{example}
我们经常不假思索地将德摩根定律用于逻辑运算符。例如，说 `$3$ 和 $7$ 都不是偶数' 相当于说 `$3$ 是奇数，$7$ 也是奇数'。前一个陈述翻译为
\[ \neg [(3 \text{为偶数}) \vee (7 \text{为偶数})]\]
后一个陈述翻译为
\[ [\neg (3 \text{为偶数})] \wedge [\neg (7 \text{为偶数})]\]
\end{example}

\begin{exercise}
\label{exNegationOfImplication}
用两种方法证明 $\neg (p \Rightarrow q) \equiv p \wedge (\neg q)$，一次使用真值表，一次使用 \Cref{exImplicationInTermsOfDisjunctionWithTruthTables} 以及德摩根定律和双重否定定律。
\end{exercise}

逻辑运算符的德摩根定律推广到有关量词的陈述，表达了 $\forall$ 和 $\exists$ 之间的对偶性，就像 $\wedge$ 和 $\vee$ 之间的对偶性一样。

\begin{theorem}[量词德摩根定律]
\label{thmDeMorganQuantifiers}
\index{de Morgan's laws!for quantifiers}
设 $p(x)$ 是一个逻辑公式，其中自由变量 $x$ 的取值范围在集合 $X$ 上。则：
\begin{enumerate}[(a)]
\item $\neg \forall x \in X,\, p(x) \equiv \exists x \in X,\, \neg p(x)$；且
\item $\neg \exists x \in X,\, p(x) \equiv \forall x \in X,\, \neg p(x)$。
\end{enumerate}
\end{theorem}

\begin{cproof}
不幸的是，由于这些逻辑公式涉及量词，我们没有可用的真值表，因此我们必须假设每个公式并导出另一个公式。

我们首先证明(b)的等价性，然后在推导出(a)结果。

\begin{itemize}
\item 假设 $\neg \exists x \in X,\, p(x)$。为了证明 $\forall x \in X,\, \neg p(x)$，给定某些 $x \in X$。如果 $p(x)$ 为真，那么我们就会有 $\exists x \in X,\, p(x)$，这与我们的主要假设相矛盾；于是我们有 $\neg p(x)$。因此 $\forall x \in X,\, \neg p(x)$ 为真。

\item 假设 $\forall x \in X,\, \neg p(x)$。为了得到矛盾，假设 $\exists x \in X,\, p(x)$ 为真。则我们有某些 $a \in X$ ，其中 $p(a)$ 为真。但是 $\neg p(a)$ 通过假设 $\forall x \in X,\, \neg p(a)$ 为真，因此我们得到了矛盾。因此 $\neg \exists x \in X,\, p(x)$ 为真。
\end{itemize}

这就证明了 $\neg \exists x \in X,\, p(x) \equiv \forall x \in X,\, \neg p(x)$。

现在，使用双重否定定律 (\Cref{thmDoubleNegation}) 从 (b) 得出 (a)：
\[ \exists x \in X,\, \neg p(x) \equiv \neg\neg \exists x \in X,\, \neg p(x) \overset{(b)}{\equiv} \neg \forall x \in X,\, \neg \neg p(x) \equiv \neg \forall x \in X,\, p(x)\]
\end{cproof}

\Cref{thmDeMorganQuantifiers}(b) 中的逻辑等价所给出的证明策略非常重要，以至于它有自己的名字。

\begin{strategy}[反例证明]
\label{strCounterexample}
\index{counterexample}
\index{proof!by counterexample}
为了证明 $\forall x \in X,\, p(x)$ 形式的命题为假，只需找到单个元素 $a \in X$ 使得 $p(a)$ 为假。元素 $a$ 被称为命题$\forall x \in X,\, p(x)$ 的\textbf{反例}。
\end{strategy}

\begin{example}
我们通过反例证明，并不是每个整数都能被质数整除。事实上，令 $x = 1$。$1$ 的唯一整数因子是 $1$ 和 $-1$，它们都不是质数，因此 $1$ 不能被任何质数整除。
\end{example}

\begin{exercise}
通过反例证明，并非所有有理数都可以表示为 $\dfrac{a}{b}$，其中 $a \in \mathbb{Z}$ 为偶数，$b \in \mathbb{Z}$ 为奇数。
\hintlabel{exNotEveryRationalIsEvenDividedByOdd}{%
找到一个有理数，其所有表示为两个整数之比的分母均为偶数。
}
\end{exercise}

我们现在已经了解了如何对逻辑运算符 $\neg$、$\wedge$、$\vee$ 和 $\Rightarrow$ 以及量词 $\forall$ 和 $\exists$ 求反。

\begin{definition}
\label{defMaximallyNegatedLogicalFormula}
\index{negation!maximal}
\index{logical formula!maximally negated}
如果唯一的否定运算符 $\neg$ 直接出现在谓词（或不涉及逻辑运算符或量词的其他命题）之前，则逻辑公式为 \textbf{最大否定}。
\end{definition}

\begin{example}
下面的命题公式为最大否定：
\[ [p \wedge (q \Rightarrow (\neg r))] \Leftrightarrow (s \wedge (\neg t))\]
事实上，所有 $\neg$ 都直接出现在命题变量之前。

然而，下面的命题公式\textit{不}是最大否定：
\[ (\neg \neg q) \Rightarrow q\]
这里，子公式 $\neg \neg q$ 在否定运算符 ($\neg \neg q$) 之前包含另一个否定运算符。然而，根据双重否定定律，这相当于 $q \Rightarrow q$，因为没有否定运算符可言，因此不存在最大否定。
\end{example}

\begin{exercise}
判断下列逻辑公式中哪些是最大否定。
\begin{enumerate}[(a)]
\item $\forall x \in X,\, (\neg p(x)) \Rightarrow \forall y \in X, \neg (r(x,y) \wedge s(x,y))$;
\item $\forall x \in X,\, (\neg p(x)) \Rightarrow \forall y \in X, (\neg r(x,y)) \vee (\neg s(x,y))$;
\item $\forall x \in \mathbb{R},\, [x > 1 \Rightarrow (\exists y \in \mathbb{R},\, [x < y \wedge \neg (x^2 \le y)])]$;
\item $\neg \exists x \in \mathbb{R},\, [x > 1 \wedge (\forall y \in \mathbb{R},\, [x < y \Rightarrow x^2 \le y])]$.
\end{enumerate}
\end{exercise}

下面的定理允许我们用最大否定逻辑公式来代替逻辑公式，这反过来又提出了证明策略，我们可以使用这些策略来证明看起来复杂的命题\textit{不成立}。

\begin{theorem}
\label{thmLogicalFormulaEquivalentToMaximallyNegated}
每个逻辑公式（仅使用我们迄今为止看到的逻辑运算符和量词构建）在逻辑上都相当于最大否定逻辑公式。
\end{theorem}

\begin{cidea}
与 \Cref{thmLogicalEquivalentIffSameTruthValues} 类似，\Cref{thmLogicalFormulaEquivalentToMaximallyNeated} 的严格证明需要某种形式的归纳论证，因此我们只给出证明思路。

到目前为止，我们看到的每个逻辑公式都是使用逻辑运算符 ${\wedge}、{\vee}、{\Rightarrow}$ 和 $\neg$ 以及量词 $\forall$ 和 $\exists$ 从谓词中构建的 --- 事实上，逻辑运算符 $\Leftrightarrow$ 是根据 $\wedge$ 和 $\Rightarrow$ 定义的，量词 $\exists$ 是根据量词 $\forall$ 和 $\exists$ 以及逻辑运算符 $\wedge$ 和 $\Rightarrow$ 定义的。

但是本节的结果允许我们将否定加入每个逻辑运算符和量词 `内部'，如下表所示。
\begin{center}
\begin{tabular}{rcll}
外部否定 & & 内部否定 & 证明 \\ \hline
$\neg (p \wedge q)$ &$\equiv$& $(\neg p) \vee (\neg q)$ & \Cref{thmDeMorganLogicalOperators}(a) \\
$\neg (p \vee q)$ &$\equiv$& $(\neg p) \wedge (\neg q)$ &  \Cref{thmDeMorganLogicalOperators}(b) \\
$\neg (p \Rightarrow q)$ &$\equiv$& $p \wedge (\neg q)$ & \Cref{exNegationOfImplication} \\
$\neg (\neg p)$ &$\equiv$& $p$ & \Cref{thmDoubleNegation} \\
$\neg \forall x \in X,\, p(x)$ &$\equiv$& $\exists x \in X,\, \neg p(x)$ & \Cref{thmDeMorganQuantifiers}(a) \\
$\neg \exists x \in X,\, p(x)$ &$\equiv$& $\forall x \in X,\, \neg p(x)$ & \Cref{thmDeMorganQuantifiers}(b)
\end{tabular}
\end{center}

重复将这些规则应用于逻辑公式最终会得到逻辑等价的、最大否定的逻辑公式。
\end{cidea}

\begin{example}
回想一下 \pageref{pConvergencePreliminary} 页中的逻辑公式，表述实数序列 $x_0, x_1, x_2, \dots$ 收敛的断言：
\[ \exists a \in \mathbb{R},\, \forall \varepsilon \in \mathbb{R},\, (\varepsilon > 0 \Rightarrow \exists N \in \mathbb{N},\, \forall n \in \mathbb{N},\, [n \ge N \Rightarrow |x_n-a| < \varepsilon])\]
我们将最大限度地否定它，以获得表达序列不收敛的逻辑公式。

让我们从头开始吧。我们一开始的否定公式是：
\[ \neg \exists a \in \mathbb{R},\, \forall \varepsilon \in \mathbb{R},\, (\varepsilon > 0 \Rightarrow \exists N \in \mathbb{N},\, \forall n \in \mathbb{N},\, [n \ge N \Rightarrow |x_n-a| < \varepsilon])\]
%
最大限度地否定逻辑公式的关键是忽略非密切相关的信息。在这里，我们要求反的表达式采用 $\neg \exists a \in \mathbb{R},\, (\text{占位内容})$ 的形式。`占位内容'是什么并不重要；重要的是我们正在否定一个存在量化的陈述，因此德摩根量词定律告诉我们，这在逻辑上等价于 $\forall a \in \mathbb{R},\, \neg (\text{占位内容})$。 我们应用这个规则并重写`占位内容'，得到：
\[ \forall a \in \mathbb{R},\, \neg \forall \varepsilon \in \mathbb{R},\, (\varepsilon > 0 \Rightarrow \exists N \in \mathbb{N},\, \forall n \in \mathbb{N},\, [n \ge N \Rightarrow |x_n-a| < \varepsilon])\]
%
接下来我们否定全称量化陈述，$\neg \forall \varepsilon \in \mathbb{R},\, (\text{占位内容})$，根据德摩根量词定律，它等价于 $\exists \varepsilon \in \mathbb{R},\, \neg (\text{占位内容})$:
\[ \forall a \in \mathbb{R},\, \exists \varepsilon \in \mathbb{R},\, \neg (\varepsilon > 0 \Rightarrow \exists N \in \mathbb{N},\, \forall n \in \mathbb{N},\, [n \ge N \Rightarrow |x_n-a| < \varepsilon])\]
%
此时，被否定的陈述的形式为 $(\text{占位A}) \Rightarrow (\text{占位B})$，它通过 \Cref{exNegationOfImplication} 否定为 $(\text{占位A}) \wedge \neg (\text{占位B})$。其中，`占位A'为 $\varepsilon > 0$，`占位B'为 $\exists N \in \mathbb{N}, \forall n \in \mathbb{N},\, [n \ge N \Rightarrow |x_n - a| < \varepsilon]$。因此执行这个否定会产生：
\[ \forall a \in \mathbb{R},\, \exists \varepsilon \in \mathbb{R},\, (\varepsilon > 0 \wedge \neg \exists N \in \mathbb{N},\, \forall n \in \mathbb{N},\, [n \ge N \Rightarrow |x_n-a| < \varepsilon])\]
%
现在我们再次否定存在量化公式，因此使用德摩根量词定律，得出：
\[ \forall a \in \mathbb{R},\, \exists \varepsilon \in \mathbb{R},\, (\varepsilon > 0 \wedge \forall N \in \mathbb{N},\, \neg \forall n \in \mathbb{N},\, [n \ge N \Rightarrow |x_n-a| < \varepsilon])\]
%
这里被否定的公式是全称量化公式，因此\textit{再次}使用德摩根量词定律，得出：
\[ \forall a \in \mathbb{R},\, \exists \varepsilon \in \mathbb{R},\, (\varepsilon > 0 \wedge \forall N \in \mathbb{N},\, \exists n \in \mathbb{N},\, \neg [n \ge N \Rightarrow |x_n-a| < \varepsilon])\]
%
我们就快成功！这里被否定的陈述是一个蕴涵式，因此应用规则 $\neg (p \Rightarrow q) \equiv p \wedge (\neg q)$ 再次产生：
\[ \forall a \in \mathbb{R},\, \exists \varepsilon \in \mathbb{R},\, (\varepsilon > 0 \wedge \forall N \in \mathbb{N},\, \exists n \in \mathbb{N},\, [n \ge N \wedge \neg (|x_n - a| < \varepsilon)])\]
%
此时，严格来说，公式已经被最大限度地否定，因为被否定的语句不涉及任何其他逻辑运算符或量词。然而，由于 $\neg (|x_n-a| < \varepsilon)$ 等价于 $|x_n - a| \ge \varepsilon$，我们可以进一步得到：
\[ \forall a \in \mathbb{R},\, \exists \varepsilon \in \mathbb{R},\, (\varepsilon > 0 \wedge \forall N \in \mathbb{N},\, \exists n \in \mathbb{N},\, [n \ge N \wedge |x_n - a| \ge \varepsilon])\]
%
这就是我们所能想到的最大限度的否定，所以我们就到此为止吧。
\end{example}

\begin{exercise}
找到一个逻辑上等价于 $\neg (p \Leftrightarrow q)$ 的最大化否定命题公式。
\hintlabel{exMaximallyNegateBiconditional}{%
首先用 $\Rightarrow$ 和 $\wedge$ 表示 $\Leftrightarrow$，如 \Cref{defBiconditional} 中所示。
}
\end{exercise}

\begin{exercise}
最大程度地否定下列逻辑公式，然后证明其为真或证明其为假。
\[ \exists x \in \mathbb{R},\, [x > 1 \wedge (\forall y \in \mathbb{R},\, [x < y \Rightarrow x^2 \le y])]\]
\end{exercise}

\subsection*{永真式}

在本章介绍的最后一个概念是\textit{永真式}，它可以被认为是矛盾的对立面。`永真'一词在口语中还有其他含义，但在符号逻辑的背景下它有一个精确定义。

\begin{definition}
\label{defTautology}
\index{tautology}
\textbf{永真式}是一个命题或逻辑公式，无论赋值给其组成命题变量和谓词的真值如何，它总为真。
\end{definition}

我们对永真式感兴趣的原因是，永真式可以出于任意原因在证明中的任意地方用作假设。

\begin{strategy}[永真假设]
设 $p$ 为一个命题。任何永真式都可以在 $p$ 的任何证明中假设。
\end{strategy}

\begin{example}
排中律 (\Cref{axLEM}) 准确地指出 $p \vee (\neg p)$ 是永真式。这意味着在证明任何结果时，我们可以根据命题是真还是假来划分情况，就像我们在 \Cref{propIfProductEvenThenSomeFactorEven} 中所做的那样。
\end{example}

\begin{example}
公式 $p \Rightarrow (q \Rightarrow p)$ 是永真式。

直接证明如下。为了证明 $p \Rightarrow (q \Rightarrow p)$ 为真，只需假设 $p$ 并导出 $q \Rightarrow p$ 即可。所以假设$p$。现在为了证明 $q \Rightarrow p$，只需假设 $q$ 并导出 $p$ 即可。所以假设$q$。但我们已经假设 $p$ 是真的！所以 $q \Rightarrow p$ 为真，因此 $p \Rightarrow (q \Rightarrow p)$ 为真。

真值表证明如下：
\begin{center}
\begin{tabular}{cc|c|c}
$p$ & $q$ & $q \Rightarrow p$ & $p \Rightarrow (q \Rightarrow p)$ \\ \hline
\TT & \TT & \TT & \TT \\
\TT & \FF & \TT & \TT \\
\FF & \TT & \FF & \TT \\
\FF & \FF & \TT & \TT
\end{tabular}
\end{center}
我们可以看到，无论 $p$ 和 $q$ 的真值如何，$p \Rightarrow (q \Rightarrow p)$ 都是真。
\end{example}

\begin{exercise}
\label{exTautologies}
证明以下每一个逻辑公式都是永真式：
\begin{enumerate}[(a)]
\item $[(p \Rightarrow q) \wedge (q \Rightarrow r)] \Rightarrow (p \Rightarrow r)$;
\item $[p \Rightarrow (q \Rightarrow r)] \Rightarrow [(p \Rightarrow q) \Rightarrow (p \Rightarrow r)]$;
\item $\exists y \in Y,\, \forall x \in X,\, p(x,y) \Rightarrow \forall x \in X,\, \exists y \in Y,\, p(x,y)$;
\item $[\neg (p \wedge q)] \Leftrightarrow [(\neg p) \vee (\neg q)]$;
\item $(\neg \forall x \in X,\, p(x)) \Leftrightarrow (\exists x \in X,\, \neg p(x))$.
\end{enumerate}
尝试解释每一个的含义，以及它如何用在证明中。
\end{exercise}

您可能已经注意到逻辑运算符和量词的德摩根定律与 \Cref{exTautologies} 的 (d) 和 (e) 之间的相似之处。它们似乎在说同样的事情，不同点在与 \Cref{exTautologies} 中我们使用 `$\Leftrightarrow$'，而在 \Cref{thmDeMorganLogicalOperators,thmDeMorganQuantifiers} 中我们使用 `$\equiv$'。不过，有一个重要的区别：如果 $p$ 和 $q$ 是逻辑公式，那么 $p \Rightarrow q$ 本身也是一个逻辑公式，我们可以将其作为一个数学对象来研究。然而，$p \equiv q$ 不是一个逻辑公式：它是一个\textit{关于}逻辑公式的断言，即逻辑公式 $p$ 和 $q$ 是等价的。

尽管如此，$\Leftrightarrow$ 和 $\equiv$ 之间存在着密切的关系 --- 这种关系总结为以下定理。

\begin{theorem}
\label{thmTautologyAndDerivation}
设 $p$ 和 $q$ 为逻辑公式。
\begin{enumerate}[(a)]
\item $q$ 可以从 $p$ 推导出当且仅当 $p \Rightarrow q$ 是永真式；
\item $p \equiv q$ 当且仅当 $p \Leftrightarrow q$ 是永真式。
\end{enumerate}
\end{theorem}

\begin{cproof}
对于 (a)，请注意，从 $p$ 推导出 $q$ 足以通过蕴涵 \introrule{\Rightarrow} 的引入规则来确定 $p \Rightarrow q$ 的真实性，因此如果 $q$ 可以从 $p$ 推导出，则 $p \Rightarrow q$ 是永真式。相反，如果 $p \Rightarrow q$ 是永真式，则 $q$ 可以使用蕴涵消除规则 \elimrule{\Rightarrow} 以及（永真式）假设 $p \Rightarrow q$ 为真从 $p$ 推导出。

接着由 (a) 推出 (b) ，因为逻辑等价是根据每个方向的推导来定义的，而 $\Leftrightarrow$ 是两个含义的合取。
\end{cproof}

%\textit{深呼吸！} 所有这些新的符号一开始可能会让人不知所措，但最终会让你收益匪浅。本章主要是教您一门新语言 --- 新符号、新术语 --- 因为没有它，我们未来的探索无法开展。如果您现在遇到困难，请不要担心：您很快就会掌握它的窍门，尤其是当我们开始在上下文中使用这种新语言时。当然，您也可以随时参考本章中的结果以供参考。

\begin{tldr}{secLogicalEquivalence}
\subsubsection*{逻辑等价}

\begin{tldrlist}
\tldritem{defLogicalEquivalence}
如果两个逻辑公式每一个都可以从另一个推导出来，则它们\textit{逻辑等价}（写作$\equiv$）。

\tldritem{defTruthTable}
命题公式的\textit{真值表}是对所有构成命题变量赋予真值下命题公式真值的列表。两个命题公式逻辑等价当且仅当它们在真值表中具有相同的列。
\end{tldrlist}

\subsubsection*{一些特定的逻辑等价}

\begin{tldrlist}
\tldritem{thmDoubleNegation}
命题的双重否定与原命题等价。这就产生了`间接'反证法：为了证明命题 $p$ 为\textit{真}，假设 $p$ 为假并推导出矛盾。

\tldritem{defContrapositive}
蕴涵式 $p \Rightarrow q$ 的逆否公式是 $(\neg q) \Rightarrow (\neg p)$。每个蕴涵式都等价于它的逆否式，因此为了证明 $p \Rightarrow q$，只需假设 $q$ 为假，并推导出 $p$ 为假。

\tldritem{thmDeMorganLogicalOperators}
德摩根逻辑运算符定律指出 $\neg (p \wedge q) \equiv (\neg p) \vee (\neg q)$ 且 $\neg (p \vee q) \equiv (\neg p) \wedge (\neg q)$。

\tldritem{exNegationOfImplication}
$p \Rightarrow q$ 的否定在逻辑上等价于 $p \wedge (\neg q)$，因此为了证明 $p \Rightarrow q$ 为假，只需证明 $p$ 为真，但 $q$ 为假。

\tldritem{thmDeMorganQuantifiers}
德摩根量词定律指出 $\neg \forall x \in X,\, p(x) \equiv \exists x \in X,\, \neg p(x)$ 且 $\neg \exists x \in X,\, p(x) \equiv \forall x \in X,\, \neg p(x)$。其中第一个给出了采用\textit{反例证明}的策略。
\end{tldrlist}

\subsubsection*{最大化否定}

\begin{tldrlist}
\tldritem{defMaximallyNegatedLogicalFormula}
如果逻辑公式不包含否定运算符（除了紧接在命题变量或谓词之前），则它被最大化否定。

\tldritem{thmLogicalFormulaEquivalentToMaximallyNegated}
本节中使用逻辑运算符和量词构建的每个逻辑公式在逻辑上都相当于最大化否定公式。
\end{tldrlist}

\subsubsection*{永真式}

\begin{tldrlist}
\tldritem{defTautology}
\textit{永真式}是一个逻辑公式，无论对其构成命题变量赋予什么真值，或者用什么值代替其自由变量，该公式都为真。永真式可以在证明中的任意地方用作假设。

\tldritem{thmTautologyAndDerivation}
命题 $p$ 可以从命题 $q$ 推导出当且仅当 $p \Rightarrow q$ 是永真式；$p \equiv q$ 当且仅当 $p \Leftrightarrow q$ 是同永真式。
\end{tldrlist}
\end{tldr}